\section{Концептуальная модель Марковского рерайтера путей (URL Rewriter)}

В рамках реализации L7-маршрутизации в балансировщике \texttt{stm-markov-proxy} утверждается логика детерминированного изменения путей запроса (URI), аналогичная механизму \texttt{rewrite} в веб-сервере Nginx.

\subsection{Принципы первой аппроксимации движка}
На данном этапе разработки архитектура рерайтера подчиняется следующим правилам:
\begin{enumerate}
    \item \textbf{Изоляция от внешних конфигураций:} Мы полностью отказываемся от парсинга внешних конфигурационных файлов (YAML/JSON) или создания кастомных DSL. Система правил \texttt{MarkovSystem} материализуется жестко в исходном коде приложения.
    \item \textbf{Отказ от регулярных выражений:} Поиск подстрок для применения правил осуществляется без использования ресурсоемких движков RegEx. Вместо этого применяется точный поиск подстрок по детерминированным алгоритмам (Кнута-Морриса-Пратта или Бойера-Мура).
    \item \textbf{Функциональная хвостовая рекурсия:} Ввиду иммутабельности данных в Haskell, ход автомата Маркова выражается через чистую хвостовую рекурсию. Функция пошагово принимает текущее состояние строки и неизменяемый список правил, возвращая новую строку на каждом успешном шаге подстановки.
\end{enumerate}

\subsection{Анализ физики памяти: Чистые структуры против мутаций}
Был детально проработан вопрос целесообразности мутации байтовых массивов на месте в целях оптимизации производительности:
\begin{itemize}
    \item \textbf{Специфика данных L7:} Длина типичного HTTP URL составляет не более 100–500 байт. При таких объемах создание новой структуры \texttt{ByteString} в куче GHC сводится к микроскопическому вызову ассемблерной команды \texttt{memcpy} (с использованием процессором инструкций SIMD). Время выполнения такой операции измеряется долями наносекунды.
    \item \textbf{Проблема изменения длины:} Мутация строки на месте порождает серьезный оверхед и риски безопасности при несовпадении длин заменяемого паттерна и подстановки (необходимость сдвига хвоста массива или вызова \texttt{realloc}). 
\end{itemize}
\textbf{Вердикт:} Использование чистых, иммутабельных преобразований \texttt{ByteString} полностью удовлетворяет требованиям Highload-систем, сохраняя абсолютную безопасность типов и избавляя рантайм от ручного управления указателями.
Use code with caution.
